\section{Background}
\label{sec:background}

\subsection{Alert Triage}
Alert triage follows detection: for each alert, the analyst determines what actually happened and whether a response is needed~\citep{nodoze}. Attack investigation reconstructs a complete attack story from a point of interest~\citep{kairos,clouseau}; triage instead ends with a single verdict, namely the cause of the alert together with the evidence that supports it. The analyst reaches the verdict by querying the environment. Queries differ in cost and in how well they separate the candidate explanations. Some are cheap and decisive for one explanation. Others, such as a threat-intelligence lookup, are slow and return the same answer for several attacks. The order of queries therefore determines how quickly an investigation ends, and the analyst must also judge when the evidence already suffices.

Consider an alert that reports a spike in authentication failures and HTTP errors on a web tier. It looks like credential stuffing, but the same symptom can come from SQL-injection probing, an authorized vulnerability scan, a health check misclassified as an attack, or an administration panel exposed to the public network. A threat-intelligence lookup on the top source may call it malicious, an answer that three of these causes share. Reading the admin panel's configuration, one cheap query, settles the last explanation at once, but only an investigator who keeps that explanation in view will run it.

\subsection{Local LLM Agents}
An LLM agent couples a language model with tools. At each step the model reads the task and the observations so far, then either calls a tool or ends the task~\citep{react}. Extensions add verbal self-reflection~\citep{reflexion} or learn when to call tools~\citep{toolformer}. In security, such agents have been applied to penetration testing~\citep{pentestgpt}, capture-the-flag challenges~\citep{cybench}, and investigation over security logs~\citep{excytin,clouseau}, mostly with large hosted models. Open-weight models such as the Qwen2.5 and Llama~3.1 families can instead run on an organization's own hardware; with 4-bit weight quantization, models of up to 72B parameters fit on a single GPU. Local deployment keeps telemetry on the premises, but it trades away model strength, and in an agent loop every step's decision rests on that weaker model.
